\section*{第 \ref{ch:b1:curves} 章——平面曲线}

\begin{solution}{exo:b1:curves:1}
$M(-t) = (t^2, -t^3)$：关于 $x$ 轴对称；只需研究 $t \geq 0$。
$x' = 2t$ 与 $y' = 3t^2$ 均 $\geq 0$：分支沿
向右上方运动，从 $(0,0)$ 延伸至无穷远。在 $t = 0$ 处速度
为零；由展开式 $x = t^2$，$y = t^3$ 可知切线为
$x$ 轴（$y/x = t \to 0$），且 $y$ 在 $x \geq 0$ 时变号：
一个朝左的\emph{\omterm{ex:b1:curves:astroid}{尖点}}。该曲线为半立方抛物线
$y^2 = x^3$。
\end{solution}

\begin{solution}{exo:b1:curves:2}
$M(t + 2\pi) = M(t) + (2\pi, 0)$：曲线重复出现，每次平移
一个车轮周长；只需研究一个周期。$x'(t) = 1 - \cos t \geq
0$，$y'(t) = \sin t$：拱形在 $\intoo{0}{\pi}$ 上升，
在 $\intoo{\pi}{2\pi}$ 上下降，于 $(\pi, 2)$ 处达到最高。奇点
满足 $x' = y' = 0$ 的点为 $t \in 2\pi\Z$，位于地面上。在 $t = 0$ 附近：
\[
x(t) = \frac{t^3}{6} + o(t^3),
\qquad
y(t) = \frac{t^2}{2} + o(t^2):
\]
$x$ 变号，$y \geq 0$，且 $\frac{x}{y^{3/2}}$ 有界：切线
为\emph{竖直}的（方向 $(0,1)$），是一个\omterm{ex:b1:curves:astroid}{尖点}，
此处被跟踪的点瞬时速度为零——这正是
无滑动滚动的物理特征。
\end{solution}

\begin{solution}{exo:b1:curves:3}
当 $t = \frac\pi4$ 时：$M = \bigl(\frac{\sqrt2}{4},
\frac{\sqrt2}{4}\bigr)$，切线方向为 $(-\cos t, \sin t) =
\frac{1}{\sqrt2}(-1, 1)$（见\cref{ex:b1:curves:astroid}）。切线
为 $y - \frac{\sqrt2}{4} = -(x - \frac{\sqrt2}{4})$，即 $x +
y = \frac{\sqrt2}{2}$。截距为 $\bigl(\frac{\sqrt2}{2},
0\bigr)$ 与 $\bigl(0, \frac{\sqrt2}{2}\bigr)$，两点间的线段
长为 $\sqrt{\frac12 + \frac12} = 1$。

（一般地，对点 $t$，$(\cos^3 t, \sin^3 t)$ 处的切线与
\omterm{prop:b1:vspaces:coordinates}{坐标}轴交于 $(\cos t, 0)$ 与 $(0, \sin t)$——请验证过这两点的直线
方向为 $(-\cos t, \sin t)$，并经过
$M(t)$——且截得的线段长为 $\sqrt{\cos^2 t +
\sin^2 t} = 1$。）
\end{solution}

\begin{solution}{exo:b1:curves:4}
$r = \cos 2\theta$：以 $\pi$ 为周期（对 $\theta$ 而言），曲线关于
两条\omterm{prop:b1:vspaces:coordinates}{坐标}轴对称；$r$ 在 $\theta = \pm\frac\pi4$ 处为零（原点处切线沿对角线方向），
且当 $\theta \in
\intoo{\frac\pi4}{\frac{3\pi}{4}}$ 时取负值，
点落在相反的射线上，于是得到四个花瓣，沿\omterm{prop:b1:vspaces:coordinates}{坐标}轴方向
即 $\theta = 0, \frac\pi2, \pi, \frac{3\pi}{2}$ 的方向，每个花瓣的最大半径为
$1$。

$r\cos\theta = 1$ 即方程 $x = 1$：\omterm{def:b1:curves:polar}{极坐标曲线} $r =
\frac{1}{\cos\theta}$ 即竖直直线 $x = 1$（恰好扫过一次，取
$\theta \in \intoo{-\frac\pi2}{\frac\pi2}$ 时）。
\end{solution}

\begin{solution}{exo:b1:curves:5}
$r(\theta) = 1 + 2\cos\theta$ 在 $\cos\theta = -\frac12$ 时为零：
$\theta = \pm\frac{2\pi}{3}$。关于 $x$ 轴对称，只需研究
$\theta \in \intcc{0}{\pi}$。当 $\theta \in
\intco{0}{\frac{2\pi}{3}}$ 时：$r > 0$，由 $3$ 递减至 $0$：
外层弧，沿与射线 $\theta =
\frac{2\pi}{3}$ 相切的方向进入原点。当 $\theta \in \intoc{\frac{2\pi}{3}}{\pi}$ 时：$r <
0$：点 $r\vec u(\theta)$ 位于相反的射线上（幅角为
$\theta - \pi \in \intoc{-\frac\pi3}{0}$），距离 $\abs r$
由 $0$ 增至 $1$：这画出一个小的\emph{内圈}
经过原点。当 $\theta = \pi$ 时 $r = -1$，此时点为
$-\vec u(\pi) = (1, 0)$，内圈在正 $x$ 轴上闭合。
草图：外层为一条类似心形的大曲线，
最远可达 $3$，在 $\theta = 0$ 处取得，其\omterm{def:b1:topology:closure}{内部}另有一个内圈
经过 $O$ 与 $(1,0)$。
\end{solution}

\begin{solution}{exo:b1:curves:6}
配方：
\[
9(x^2 - 4x) + 25(y^2 - 2y) = 164
\iff 9(x-2)^2 + 25(y-1)^2 = 164 + 36 + 25 = 225 .
\]
两边除以 $225$，得 $\frac{(x-2)^2}{25} + \frac{(y-1)^2}{9} = 1$，表示一个
中心为 $(2, 1)$ 的椭圆，半轴 $a = 5$（水平方向），$b =
3$；$c = \sqrt{25 - 9} = 4$，焦点为 $(2 \pm 4,\, 1) = (-2, 1)$ 与
$(6, 1)$；离心率 $e = \frac45$。
\end{solution}

\begin{solution}{exo:b1:curves:7}
焦点位于原点，准线 $D: x = d$，其中 $d = \frac pe > 0$。
设点 $M = (r\cos\theta, r\sin\theta)$，$r > 0$：
\[
MF = r,
\qquad
d(M, D) = \abs{d - r\cos\theta},
\]
而\omterm{def:b1:curves:conics}{圆锥曲线}条件 $MF = e\,d(M, D)$，在 $M$
位于焦点一侧（$r\cos\theta < d$）的情形下，可化为 $r
= e(d - r\cos\theta)$，即
\[
r(1 + e\cos\theta) = ed = p,
\qquad
r = \frac{p}{1 + e\cos\theta} .
\]
当 $e < 1$ 时分母永不为零，$\theta$ 可取全体
（椭圆）。当 $e = 1$ 时：$\theta \neq \pi$（抛物线，开口朝向
准线的远侧）。当 $e > 1$ 时需要 $1 + e\cos\theta > 0$，
即 $\theta \in \intoo{-\theta_0}{\theta_0}$，其中 $\theta_0 =
\arccos\bigl(-\frac1e\bigr)$：给出双曲线的一个\emph{分支}
（另一分支对应于 $r < 0$ 的符号选取，或对应于
第二个焦点）。
\end{solution}

\begin{solution}{exo:b1:curves:8}
\emph{园丁：}将长为 $2a$ 的绳子两端分别系在两根木桩 $F,
F'$ 上，用笔尖 $M$ 将绳子拉紧：约束条件恰为
$MF + MF' = 2a$，因此所画曲线为椭圆。

\emph{反射性质：}设 $t \mapsto M(t)$ 为正则
参数化，并记 $u(t) = \frac{M(t) - F}{\norm{M(t) - F}}$，
$u'(t)$ 是对 $F'$ 类似定义的单位向量。对
$\norm{M - F} + \norm{M - F'} = 2a$ 两边求导，并利用 $\frac{\dd}{\dd
t}\norm{M - F} = \bigl\langle \frac{M - F}{\norm{M-F}},\,
M'\bigr\rangle$（对 $\sqrt{\langle\cdot,\cdot\rangle}$ 用链式法则）：
\[
\bigl\langle u + u',\, M'\bigr\rangle = 0 .
\]
于是切方向 $M'$ 正交于两焦半径的
平分线方向 $u + u'$：切线
与 $MF$、$MF'$ 所成角相等。（从一个焦点发出的光线
经椭圆反射后射向另一焦点。）
\end{solution}

\begin{solution}{exo:b1:curves:9}
\begin{enumerate}
  \item 定义域：$\cos 2\theta \geq 0$，即 $\theta \in
        \intcc{-\frac\pi4}{\frac\pi4} \cup
        \intcc{\frac{3\pi}{4}}{\frac{5\pi}{4}}$。对称性：$\theta
        \mapsto -\theta$（关于 $x$ 轴）及 $\theta \mapsto \pi -
        \theta$（关于 $y$ 轴）：只需研究 $\intcc{0}{\frac\pi4}$；$r$
        由 $1$ 减小至 $0$。当 $\theta = \pm\frac\pi4$ 时：$r =
        0$，原点处切线沿两条对角线方向。曲线
        形如 $\infty$ 符号：两个对称的环在 $O$ 处相交，
        并经过点 $(\pm 1, 0)$。
  \item 设 $F, F' = (\pm c, 0)$，$c = \frac{1}{\sqrt 2}$，点 $M
        = (r\cos\theta, r\sin\theta)$:
        \[
        MF^2\, MF'^2
        = \bigl((r\cos\theta - c)^2 + r^2\sin^2\theta\bigr)
        \bigl((r\cos\theta + c)^2 + r^2\sin^2\theta\bigr)
        = (r^2 + c^2)^2 - (2rc\cos\theta)^2 ,
        \]
        由恒等式 $(A - B)(A + B) = A^2 - B^2$，取 $A = r^2
        + c^2$，$B = 2rc\cos\theta$。代入 $c^2 = \frac12$：
        \[
        MF^2 MF'^2 = \Bigl(r^2 + \frac12\Bigr)^2 - 2r^2\cos^2\theta
        = r^4 + r^2\bigl(1 - 2\cos^2\theta\bigr) + \frac14
        = r^4 - r^2\cos 2\theta + \frac14 .
        \]
        在双纽线 $r^2 = \cos 2\theta$ 上，前两项
        相互抵消，剩下 $MF^2MF'^2 = \frac14$，即\ $MF \cdot MF' =
        \frac12$。反过来，将上述计算逆推可知
        轨迹方程 $MF\,MF' = \frac12$ 等价于 $r^2 = \cos
        2\theta$（要求 $r \neq 0$；$O$ 两式均满足）。
\end{enumerate}
\end{solution}

\begin{solution}{exo:b1:curves:10}
$r = 1 + \cos\theta$, $r' = -\sin\theta$:
\[
r^2 + r'^2 = 1 + 2\cos\theta + \cos^2\theta + \sin^2\theta
= 2 + 2\cos\theta = 4\cos^2\frac\theta2 ,
\]
故 $\sqrt{r^2 + r'^2} = 2\abs{\cos\frac\theta2}$，由
关于 $x$ 轴的对称性，
\[
L = \int_0^{2\pi} 2\Bigl|\cos\frac\theta2\Bigr|\,\dd\theta
= 2\int_0^{\pi} 2\cos\frac\theta2\,\dd\theta
= 2\Bigl[4\sin\frac\theta2\Bigr]_0^{\pi} = 8 .
\]
另一个不含 $\pi$ 的代数周长。
\end{solution}

\begin{solution}{exo:b1:curves:11}
点 $(a\cos t, b\sin t)$ 满足方程：
$\cos^2 t + \sin^2 t = 1$。该直线的法向量为
$\bigl(\frac{\cos t}a, \frac{\sin t}b\bigr)$，它与
速度 $(-a\sin t, b\cos t)$ 的\omterm{def:b1:euclid:def}{内积}为 $-\sin t\cos t +
\sin t\cos t = 0$，故该直线过此点，
且方向为切方向——即它就是切线。设 $(x_0, y_0)$ 为椭圆上的
点，记 $\cos t = \frac{x_0}a$，$\sin t = \frac{y_0}b$，
代入得：
\[
\frac{x\,x_0}{a^2} + \frac{y\,y_0}{b^2} = 1 ,
\]
由椭圆方程经“拆分”得到：$x^2 \mapsto
x\,x_0$ 及 $y^2 \mapsto y\,y_0$。
\end{solution}

\begin{solution}{exo:b1:curves:12}
\begin{enumerate}
  \item $r' = k\eu^{k\theta}$，故 $\tan V = \frac{r}{r'} =
        \frac1k$，为常数。该螺线与从原点出发的每条射线
        相交于同一夹角 $V = \arctan\frac1k$。
  \item 用复数记号，螺线为 $\{\eu^{k\theta}
        \eu^{\iu\theta} : \theta \in \R\}$。旋转 $c$
        相当于乘以 $\eu^{\iu c}$：
        \[
        \eu^{k\theta}\eu^{\iu(\theta + c)}
        = \eu^{-kc}\,\eu^{k(\theta + c)}\eu^{\iu(\theta+c)} ,
        \]
        当 $\theta + c$ 取遍 $\R$ 时，这描述的是
        $\eu^{-kc}$ 倍的螺线：旋转 $=$ 伸缩。
        除直线与圆外，没有其他光滑曲线
        具有此性质。
  \item $\sqrt{r^2 + r'^2} = \sqrt{1 + k^2}\,\eu^{k\theta}$，故
        在 $\intcc{A}{\theta_0}$ 上的\omterm{def:b1:curves:arclength}{弧长}为
        $\frac{\sqrt{1+k^2}}{k}\bigl(\eu^{k\theta_0} -
        \eu^{kA}\bigr)$，当 $A \to -\infty$ 时：
        \[
        L = \frac{\sqrt{1 + k^2}}{k}\,\eu^{k\theta_0} < \infty :
        \]
        绕原点有无穷多圈，总长却有限
        （各圈按几何\omterm{def:b1:series:def}{级数}递减）。
\end{enumerate}
\end{solution}

\begin{solution}{pb:b1:curves:1}
\textbf{1.} 纯滚动是指接触点所经过的距离
恰等于展开的轮缘\omterm{def:b1:curves:arclength}{弧长}：转过 $t$ 后，
圆心位于 $(Rt, R)$。标记点位于轮缘上，
与竖直向下方向成角 $t$，位于其\emph{后方}
（车轮前进时顺时针转动）：
\[
M(t) = \Omega(t) + R(-\sin t, -\cos t)
= \bigl(R(t - \sin t),\ R(1 - \cos t)\bigr),
\]
此式在 $t = 0$（$M = (0,0)$）及 $t = \pi$（$M =
(\pi R, 2R)$，即最高点）处均成立。

\textbf{2.} $f'(t) = R(1 - \cos t,\ \sin t)$，且
\[
\norm{f'}^2 = R^2\bigl((1-\cos t)^2 + \sin^2 t\bigr)
= 2R^2(1 - \cos t) = 4R^2\sin^2\frac t2 ,
\]
于是 $\norm{f'} = 2R\abs{\sin\frac t2}$。奇点位于 $t \in
2\pi\Z$ 处，即地面上的\omterm{ex:b1:curves:astroid}{尖点}（见\cref{exo:b1:curves:2}）；
拱顶对应 $t = \pi$，此处速率 $2R$ 最大。

\textbf{3.} 半角分解：
\[
f'(t) = 2R\sin\frac t2\,\Bigl(\sin\frac t2,\ \cos\frac t2\Bigr),
\quad
M - C = 2R\sin\frac t2\,\Bigl(-\cos\frac t2,\ \sin\frac t2\Bigr),
\]
\[
T - M = \bigl(R\sin t,\ R(1 + \cos t)\bigr)
= 2R\cos\frac t2\,\Bigl(\sin\frac t2,\ \cos\frac t2\Bigr).
\]
当 $0 < t < 2\pi$ 时 $\sin\frac t2 \neq 0$，$M - C$ 与
切线方向 $\bigl(\sin\frac t2, \cos\frac
t2\bigr)$ 正交（其\omterm{def:b1:euclid:def}{内积}为 $-\sin\frac t2\cos\frac t2 +
\sin\frac t2\cos\frac t2 = 0$），而 $T - M$ 与之平行。
连接轮顶的弦\emph{是}切线；
连接接触点的弦则是法线。

\textbf{4.} 每一时刻，车轮绕其接触点转动
（该点速度为零：纯滚动），
故车轮上每一刚性点的运动方向都与其
到 $C$ 的连线正交，且速率（角速度为 $1$）恰等于
该距离。验证：$\norm{M - C} = 2R\abs{\sin\frac t2} =
\norm{f'(t)}$。

\textbf{5.} $2R\,y(t) = 2R\cdot 2R\sin^2\frac t2 =
4R^2\sin^2\frac t2 = \norm{f'(t)}^2$。高度为 $y$ 处的速率为
$\sqrt{2R\,y}$——形式上即自由落体的 $v = \sqrt{2gh}$ 规律，
只是此处以地面充当顶板；第四部分将这一观察
化为严格的力学理论。

\textbf{6.} 由 \cref{def:b1:curves:arclength} 及第 2 小题
（$\sin\frac t2 \geq 0$ 在 $\intcc{0}{2\pi}$ 上成立）：
\[
L = \int_0^{2\pi} 2R\sin\frac t2\,\dd t
= 2R\Bigl[-2\cos\frac t2\Bigr]_0^{2\pi} = 2R(2 + 2) = 8R .
\]
雷恩定理：恰为四个直径之长。

\textbf{7.} $s(t) = \int_0^t 2R\sin\frac u2\,\dd u =
4R\bigl(1 - \cos\frac t2\bigr)$; $s(2\pi) = 4R(1+1) = 8R$,
结果一致。

\textbf{8.} $\sigma(t) = \abs{s(t) - s(\pi)} = \abs{4R(1 -
\cos\frac t2) - 4R} = 4R\abs{\cos\frac t2}$，故 $\sigma^2 =
16R^2\cos^2\frac t2$；且 $2R - y = R(1 + \cos t) =
2R\cos^2\frac t2$，于是 $8R(2R - y) = 16R^2\cos^2\frac t2 =
\sigma^2$。

\textbf{9.} 在\omterm{ex:b1:curves:astroid}{尖点} $t = 0$（或 $2\pi$）处 $\sigma = 4R$，恰为
$8R$ 之半：拱顶将拱弧平分。两种计算中，速率
为 $\abs{\text{ 三角多项式，自变量为 } t/2}$，其原函数
仍是三角函数：长度是余弦的
\emph{值}之差，为 $R$ 的有理数倍，
无需度量圆弧，故不含 $\pi$。
圆周本身速率恒定，故其长度\omterm{thm:b1:integration:def}{积分}
给出整个\omterm{prop:b1:reals:intervals}{区间}的长度 $2\pi$；摆线的速率
在端点处为零，其\omterm{thm:b1:integration:def}{积分}可代数地求得。

\textbf{10.} 拱弧扫过时 $x$ 由 $0$ 增至
$2\pi R$（$x' = R(1 - \cos t) \geq 0$，仅在
孤立点处为零）。将 $x = x(t)$ 代入面积\omterm{thm:b1:integration:def}{积分}
$\int_0^{2\pi R} y\,\dd x$（见\cref{thm:b1:integration:parts}）
得 $A = \int_0^{2\pi} y(t)\,x'(t)\,\dd t$。

\textbf{11.}
\[
A = \int_0^{2\pi} R(1 - \cos t)\cdot R(1 - \cos t)\,\dd t
= R^2\int_0^{2\pi}\bigl(1 - 2\cos t + \cos^2 t\bigr)\dd t
= R^2\Bigl(2\pi - 0 + \pi\Bigr) = 3\pi R^2 ,
\]
利用 $\int_0^{2\pi}\cos^2 = \pi$。恰为三个车轮圆面积：
伽利略的称重法称其“约为 $3$”；罗贝瓦尔的计算
则断言“恰好”。

\textbf{12.} 取 $x = \cos^3 t$，$y = \sin^3 t$，则 $y\,x' =
-3\sin^4 t\cos^2 t$。积化和差：
\[
\sin^4 t\cos^2 t = \frac{1 - \cos 2t}{2}\cdot\frac{\sin^2
2t}{4} = \frac{\sin^2 2t}{8} - \frac{\sin^2 2t\cos 2t}{8} ,
\]
在 $\intcc{0}{2\pi}$ 上：$\int\sin^2 2t = \pi$，$\int \sin^2
2t\cos 2t = \bigl[\frac{\sin^3 2t}{6}\bigr] = 0$。故 $\oint
y\,\dd x = -3\cdot\frac\pi8$，所围面积为
$\frac{3\pi}8$（符号体现逆时针
方向）。

\textbf{13.} 对 $(\cos t, \sin t)$，令 $t$ 由 $\pi$ 变到 $0$，
$x$ 自 $-1$ 增至 $1$，且
\[
\int_\pi^0 \sin t\cdot(-\sin t)\,\dd t
= \int_0^\pi \sin^2 t\,\dd t = \frac\pi2 :
\]
公式恰给出上半圆盘的面积，理应如此。

\textbf{14.} $x' = R(1 + \cos t) = 2R\cos^2\frac t2$, $y' =
R\sin t = 2R\sin\frac t2\cos\frac t2$，故 $\norm{f'} =
2R\cos\frac t2$（在 $\intcc{-\pi}{\pi}$ 上非负）。\omterm{def:b1:curves:arclength}{弧长}自
最低点起算：$s(t) = \int_0^t 2R\cos\frac u2\,\dd u =
4R\sin\frac t2$。于是
\[
y = R(1 - \cos t) = 2R\sin^2\frac t2
= 2R\Bigl(\frac{s}{4R}\Bigr)^{\!2} = \frac{s^2}{8R} .
\]

\textbf{15.} 能量守恒，取 $v = \frac{\dd s}{\dd\tau}$
及 $y = \frac{s^2}{8R}$，$y_0 = \frac{s_0^2}{8R}$：
\[
\Bigl(\frac{\dd s}{\dd\tau}\Bigr)^{\!2} = 2g(y_0 - y)
= \frac{2g}{8R}\bigl(s_0^2 - s^2\bigr)
= \frac{g}{4R}\bigl(s_0^2 - s^2\bigr).
\]

\textbf{16.} 对 $\tau$ 求导：$2s's'' =
-\frac{g}{4R}\,2ss'$，故凡 $s' \neq 0$ 之处（由\omterm{def:b1:continuity:continuous}{连续}性知
处处如此），$s'' = -\frac{g}{4R}s$：即调和振子。
由 \cref{thm:b1:diffeq:homogeneous2} 可知，$s(\tau) =
A\cos\omega\tau + B\sin\omega\tau$，其中 $\omega =
\sqrt{g/(4R)}$；初始条件 $s(0) = s_0$，$s'(0) = 0$
给出 $s(\tau) = s_0\cos\omega\tau$。

\textbf{17.} 小珠到达最低点时 $s = 0$，即
$\omega\tau = \frac\pi2$ 时：
\[
T_{\downarrow} = \frac{\pi}{2\omega}
= \frac\pi2\sqrt{\frac{4R}{g}} = \pi\sqrt{\frac Rg}\,,
\]
与 $s_0$ 无关：从碗中任一处释放，所有
小珠同时到达——即等时曲线。

\textbf{18.} 代入 $s = s_0\cos\omega\tau$：左边为
$s_0^2\omega^2\sin^2\omega\tau$，右边为
$\frac{g}{4R}s_0^2(1 - \cos^2\omega\tau) =
s_0^2\omega^2\sin^2\omega\tau$：恰好相等，故未引入
增根。对圆弧而言，$y$ 并不
与 $s^2$ 成正比（$y = \ell(1 - \cos\frac s\ell) =
\frac{s^2}{2\ell} - \frac{s^4}{24\ell^3} + \dots$）：
回复项仅\emph{近似}线性，故
圆摆的周期随振幅而漂移，而
摆线的周期则严格恒定。

\textbf{19.} 由 $\pi\sqrt{R/g} = 1$ 得 $R = \frac{g}{\pi^2}
\approx \frac{9.81}{9.87} \approx 0.994$ 米——几乎恰为
一米。一次完整振动所需时间为 $\frac{2\pi}\omega =
2\pi\sqrt{4R/g}$，恰为小角度公式
$2\pi\sqrt{\ell/g}$（对应摆长 $\ell = 4R$ 的单摆）：
惠更斯将其单摆悬挂于比例恰为此值的
摆线形颊板之上。

\textbf{20.} 当 $t = \frac\pi2$，$R = 1$ 时：$M = \bigl(\frac\pi2 -
1,\ 1\bigr)$，$T = \bigl(\frac\pi2,\ 2\bigr)$，故 $T - M = (1,
1)$。又 $f'(\frac\pi2) = (1 - 0,\ 1) = (1, 1)$：指向
顶点的弦即为速度，正如所言。

\textbf{21.} 对于广义旋轮线，$x'(t) = R - d\cos t$ 及 $y'(t)
= d\sin t$。若 $d < R$，则 $x' \geq R - d > 0$，动点
始终向前推进；速度永不为零（其第一分量
为正）：无奇点，是正规波。若 $d > R$，则 $x'(0)
= R - d < 0 < R + d = x'(\pi)$，故动点在
接触时刻附近向后运动，其余时刻向前运动，曲线
自相交而成环。铁路车轮轮缘上的一点
（位于轨顶之下，$d > R$）每转一圈
都有一段倒行。

\textbf{22.} 自\omterm{ex:b1:curves:astroid}{尖点} $(\pi R, 2R)$ 到原点的弦
长为 $L = R\sqrt{\pi^2 + 4}$，倾角 $\alpha$ 满足
$\sin\alpha = \frac{2R}{L}$。小珠自静止沿弦滑下，加速度为常数
$g\sin\alpha$：$L = \frac12 g\sin\alpha\,T^2$，于是
\[
T = \sqrt{\frac{2L}{g\sin\alpha}} = \sqrt{\frac{2L^2}{2Rg}}
= \frac{L}{\sqrt{Rg}} = \sqrt{\pi^2 + 4}\,\sqrt{\frac Rg}
\approx 3.72\sqrt{\frac Rg}\,,
\]
与摆线的 $\pi\sqrt{R/g} \approx 3.14\sqrt{R/g}$ 相比：
弯曲路径更快。事实上，它是所有可能路径中最快的
——最速降线——，是变分法中的一个定理，
超出本分册范围。

\textbf{23.} $\dfrac{\dd y}{\dd x} = \dfrac{y'}{x'} =
\dfrac{\sin t}{1 - \cos t} = \cot\dfrac t2$（半角
公式）。于是
\[
\frac{\dd^2y}{\dd x^2}
= \frac{\dd}{\dd t}\Bigl(\cot\frac t2\Bigr)\cdot\frac{1}{x'(t)}
= -\frac{1}{2\sin^2\frac t2}\cdot\frac{1}{2R\sin^2\frac t2}
= -\frac{1}{4R\sin^4\frac t2} < 0 :
\]
整个拱上为凹；当 $t \to 0^+$ 或 $2\pi^-$ 时，
斜率 $\cot\frac t2 \to \pm\infty$，
即\omterm{ex:b1:curves:astroid}{尖点}处切线竖直。

\textbf{24.} 弦的方向为 $T - M \parallel
\bigl(\sin\frac t2, \cos\frac t2\bigr)$，其趋向 $(0, 1)$
当 $t \to 0^{+}$ 时，\omterm{ex:b1:curves:astroid}{尖点}处切线竖直——
由纯几何即可得出，无需泰勒展开。

\textbf{25.} (i) 雷恩定理，$L = 8R$；罗贝瓦尔面积，$A =
3\pi R^2$（伽利略所猜测的比值为 $3$）；惠更斯
等时曲线，$T_\downarrow = \pi\sqrt{R/g}$。(ii) 所有计算
均基于半角分解 $f' =
2R\sin\frac t2(\sin\frac t2, \cos\frac t2)$ 及其碗形类比
——同一恒等式同时支撑了切线、\omterm{def:b1:curves:arclength}{弧长}与
下落时间的计算。(iii) 内蕴关系 $y = s^2/8R$ 将能量方程
化为 $s'' = -\frac{g}{4R}s$，即一个常系数线性方程，
其解恰为
等时的。(iv) 第一部分用到
\cref{ch:b1:derivative}中的微分学，第二、三部分用到
\cref{ch:b1:integration}中的\omterm{thm:b1:integration:def}{积分}，第四部分用到
\cref{ch:b1:diffeq}中的微分方程——摆线即本书
课程内容汇聚于一条曲线之中。
\end{solution}
